%%%%%%%%%%%%%%%%%%%%%%%%%%%%%%%%%
%
%       EXERCÍCIO
%
%%%%%%%%%%%%%%%%%%%%%%%%%%%%%%%%%

\ifdefstring{\atividade}{prova}{%
    \ifdefstring{\modo}{objetivo}{%
        \renewcommand{\valorquestao}{\ValorQObj\ ponto}
    }{%
        \renewcommand{\valorquestao}{\ValorQDisc\ pontos}
    }%
}%

\begin{exercicioBanco}[\valorquestao]
Durante um período de 12 anos, o valor de revenda de um veículo decresce linearmente com o tempo, isto é, o valor \(y\), em função do tempo \(x\), é dado por uma função afim \(y = ax + b\). Sabe-se que, daqui a 3 anos, o valor do veículo será R\$ \(18000{,}00\), e daqui a 9 anos será R\$ \(9000{,}00\). Determine qual será o valor do veículo daqui a 6 anos.

% Define as alternativas
\newcommand{\alternativas}{%
    \begin{center}
        \begin{tabularx}{\textwidth}{XXXXX}
            (a) R\$ \(11000{,}00\). &
            (b) R\$ \(12000{,}00\). &
            (c) R\$ \(13000{,}00\). &
            (d) R\$ \(14000{,}00\). &
            (e) R\$ \(15000{,}00\).
        \end{tabularx}
    \end{center}
}

% Resposta correta
\newcommand{\resposta}{C}

% Lógica condicional para exibição
\ifdefstring{\atividade}{lista}{%
    \alternativas
    \vspace{0.5em}
    
    \noindent\textbf{Resposta:} letra \textbf{\resposta}.
}{%
    \ifdefstring{\modo}{objetivo}{%
        \alternativas
    }{%
        % modo = discursiva → não mostra alternativas
    }
}
\end{exercicioBanco}

